\chapter{Strategy observables and the expanded space}
\label{chap:qpbt-observables}
% This chapter covers the first three subsections of the source appendix
% ``Analysis of the Pauli basis test''
% (references/qpbt-paper/14_analysis_of_the_pauli_basis_test.tex, lines 1-679).
% The paper's subsection labels sec:qld-prelim, sec:commutation and
% sec:expanding are kept on the sections below for cross-reference fidelity.
% The source restates thm:pauli verbatim at the head of the appendix under the
% label thm:pauli-appendix (lines 6-20); the statement is the same theorem, and
% all references here are to thm:pauli as stated in Chapter chap:qpbt-test.

This chapter and the two that follow contain the proof of the soundness
theorem for the Pauli basis test, Theorem~\ref{thm:pauli}, stated in
Chapter~\ref{chap:qpbt-test}. Starting from a strategy that succeeds in the
game $\mathfrak{G}^{\mathrm{Pauli}}_{(q,m,d)}$ with high probability, the proof
constructs a representation of the Pauli group on the players' Hilbert space
in several stages. In Section~\ref{sec:commutation} the strategy measurements
are used to define a set of observables that correspond to a subset of the
Pauli $X$ and $Z$ observables, and these observables are shown to
approximately satisfy the associated Pauli group relations. In
Section~\ref{sec:expanding}, ancillary registers containing EPR states are
adjoined to the players' state, and a new set of approximately commuting
observables acting on the expanded state is defined.
Chapter~\ref{chap:qpbt-combining} combines the approximately commuting
observables into a strategy for the \emph{classical} low individual degree
test and applies the soundness theorem of that test to construct a single
projective measurement that simultaneously measures all of them.
Chapter~\ref{chap:qpbt-extraction} uses this simultaneous measurement to
construct an exact representation of the Pauli group, shows that it is close
to the original strategy observables on the expanded state, and deduces
Theorem~\ref{thm:pauli}. Section~\ref{sec:qld-prelim} below collects
preliminary definitions and lemmas used throughout the analysis.

\section{Preliminaries}
\label{sec:qld-prelim}

\begin{definition}[Individual and total degree polynomial classes]
	\label{def:ideg-deg-polynomials}
	\uses{def:line, def:line-representative, def:polyfunc}
	Given integers $d, m \geq 0$ and a prime power $q$, denote by
	$\mathrm{ideg}_{d,m}(\F_q)$ the set of polynomials in $m$ variables over
	$\F_q$ with \emph{individual} degree at most $d$ in each variable. When
	the field is clear from context we write $\mathrm{ideg}_{d,m}$, and when
	the number of variables is also clear from context, $\mathrm{ideg}_d$. We
	denote by $\deg_d$ the set of polynomials with \emph{total} degree at
	most $d$.
	% The line-restricted classes below are introduced only implicitly in the
	% source, in the statement of the lemma labelled lem:qld-comm-line-cons,
	% which reads them as polynomials "defined on the line". The coefficient
	% reading adopted here follows the answer format of the decision procedure
	% of the low individual degree game (source figure fig:ld-decider) and the
	% representative convention of def:polyfunc; see the remark following this
	% definition.
	For a line $\ell \subseteq \F_q^m$ (Definition~\ref{def:line}) presented
	by a base point $u_0$ and a direction $v$ --- for the lines arising as
	question contents, the pair read off the line's description, whose base
	point is the canonical representative of
	Definition~\ref{def:line-representative} --- and an integer $c \geq 0$,
	we write $\deg_c(\ell)$ for the set of univariate polynomials of degree
	at most $c$, each given by its list of $c+1$ coefficients and identified
	by its polynomial representative rather than by the function it induces,
	as in Definition~\ref{def:polyfunc}; the argument $\ell$ records the
	line to which such an answer refers, through the parameterization
	$t \mapsto u_0 + tv$. The cases used below are $\deg_d(\ell)$ and
	$\deg_{md}(\ell)$; whenever $d \leq md$, extending coefficient lists by
	zeros gives the inclusion $\deg_d(\ell) \subseteq \deg_{md}(\ell)$. For
	$f \in \deg_c(\ell)$, $u \in \ell$, and $a \in \F_q$, we say that $f$
	\emph{evaluates to $a$ at $u$}, written $f(u) = a$, if $f(t) = a$ for
	every $t \in \F_q$ with $u = u_0 + tv$. When $v \neq 0$ this $t$ is
	unique, so every $f$ evaluates at every point of $\ell$; when $v = 0$,
	so that $\ell = \{u_0\}$, exactly those $f$ that induce a constant
	function on $\F_q$ evaluate at $u_0$, to their constant value. An
	operator indexed by $f(u)$ is understood to be the zero operator when
	$f$ does not evaluate at $u$, and for a family $\{N_f\}$ indexed by
	$f \in \deg_c(\ell)$ we write
	$N_{[\mathrm{eval}_u(\cdot)=a]} = \sum_{f \colon f(u) = a} N_f$; when
	$v = 0$ the classes indexed by $a \in \F_q$ form a sub-measurement
	rather than a partition of the outcomes. The \emph{completed} family of
	evaluation classes is the family indexed by $a \in \F_q \cup \{\bot\}$
	obtained by adjoining the additional class
	$N_{[\mathrm{eval}_u(\cdot)=\bot]} = \sum_{f} N_f$, the sum over those
	$f$ that do not evaluate at $u$; when $\{N_f\}$ is a POVM,
	respectively a projective measurement, the completed classes form one
	as well, and when $v \neq 0$ the class
	$N_{[\mathrm{eval}_u(\cdot)=\bot]}$ is the zero operator.
\end{definition}

\begin{remark}[Line answers as coefficient lists]
	\label{rem:deg-line-representatives}
	The source refers to the elements of $\deg_{md}(\ell)$ as polynomials
	``defined on the line $\ell$'', that is, as functions on the point set
	of $\ell$, while its decision procedure expects each line answer as a
	list of coefficients. The two readings differ. When $md \geq q$,
	distinct coefficient lists induce the same function $\F_q \to \F_q$;
	and when $\ell$ has direction $v = 0$ --- a case the question
	distributions produce with positive probability
	(Remark~\ref{rem:ld-win-zero-direction}) --- the point set of $\ell$ is
	the singleton $\{u_0\}$, so a function on $\ell$ retains only a single
	field element of the answer. Definition~\ref{def:ideg-deg-polynomials}
	therefore adopts the coefficient reading, under which the strategy's
	line measurements are indexed by exactly the answers of the prescribed
	form, and its evaluation convention reproduces the universal
	quantification of the line-versus-point clauses of the win predicates
	(Definition~\ref{def:ld-win-predicate}).
\end{remark}

Throughout this chapter $q = 2^k$ with $k$ odd is an admissible field size
(Definition~\ref{def:admissible-size}) and $M = 2^m$ is the number of
$q$-ary qudits tested for per player, as in
Chapters~\ref{chap:qpbt-test} and~\ref{chap:qpbt-extraction}. (The letter
$M$ also denotes the measurement operators of the strategy under analysis,
as in the tuple $\mathcal{S} = (\ket{\psi}, M)$ below and in the
super- and subscripted individual operators such as
$M^{(\mathrm{Point},W),u}_a$; the meaning is determined by context, as in
the source.) Recall from
Chapter~\ref{chap:qpbt-algebra} the low-degree code
(see Definition~\ref{def:low-degree-encoding}), which maps each vector
$h \in \F_q^{M}$, with
coordinates indexed by binary strings of length $m$, to a polynomial
$g_h \in \mathrm{ideg}_{d,m}(\F_q)$ such that $g_h(u) = h_u$ for every
$u \in \{0,1\}^m$; for $x \in \F_q^m$ the associated vector
$\mathrm{ind}_m(x) \in \F_q^{M}$ satisfies
$g_h(x) = h \cdot \mathrm{ind}_m(x)$ for all $h$. Recall also the finite
field trace $\tr = \tr_{q \to 2} \colon \F_q \to \F_2$
(see Definition~\ref{def:subfield-trace}).

\begin{definition}[Commuting and anticommuting tuples]
	\label{def:anticommuting-tuple}
	% The quantity gamma is introduced in the source at eq:gamma-value, inside
	% the completeness analysis of the Pauli basis test; the definition below
	% is self-contained and does not depend on that occurrence.
	\uses{def:low-degree-encoding, def:subfield-trace}
	For a tuple
	$\omega = (u_X, u_Z, r_X, r_Z) \in (\F_q^m)^2 \times (\F_q)^2$, let
	\[
		\gamma(\omega) \,=\, \tr\big( (r_X \cdot \mathrm{ind}_m(u_X)) \cdot (r_Z \cdot \mathrm{ind}_m(u_Z)) \big)\,,
	\]
	where $r_W \cdot \mathrm{ind}_m(u_W)$ is the scalar multiple of the
	vector $\mathrm{ind}_m(u_W) \in \F_q^{M}$, the middle $\cdot$ is the
	$\F_q$-inner product of two vectors in $\F_q^{M}$, and
	$\tr = \tr_{q \to 2}$. The tuple $\omega$ is called \emph{anticommuting}
	if $\gamma(\omega) \neq 0$, and \emph{commuting} if $\gamma(\omega) = 0$.
\end{definition}

The sign $(-1)^{\gamma(\omega)}$ is exactly the sign with which the
generalized Pauli observables $\tau^X(r_X \cdot \mathrm{ind}_m(u_X))$ and
$\tau^Z(r_Z \cdot \mathrm{ind}_m(u_Z))$ commute
(see Lemma~\ref{lem:twisted-commutation});
this is the source of the terminology.

\begin{lemma}[Probability of anticommuting tuples]
	\label{fact:omega-anticomm-prob}
	% The source states this result in a `fact' environment; it is stated here
	% as a lemma, and the paper label is kept verbatim.
	% The exact value of the anticommuting probability and the constant lower
	% bounds under the hypothesis m <= q are additions relative to the source,
	% which asserts only the two bounds displayed at the end of the statement;
	% the source states those bounds without the hypothesis m, d >= 1, under
	% which alone they hold (see the remark following this lemma). The
	% unconditional constant bounds are what the proof of the lemma labelled
	% lem:qld-win-implications uses.
	\uses{def:anticommuting-tuple}
	A tuple $\omega = (u_X,u_Z,r_X,r_Z)$ sampled uniformly at random from
	$(\F_q^m)^2 \times (\F_q)^2$ is anticommuting with probability exactly
	\[
		\big(1 - q^{-1}\big)^{m+1} \cdot \frac{1}{2}\,,
	\]
	and commuting with the complementary probability, which is at least
	$\frac12$. In particular, if $m \leq q$, then $\omega$ is anticommuting
	with probability at least $2^{-4}$ and commuting with probability at
	least $\frac12$. If $m, d \geq 1$, then the probability that $\omega$ is
	anticommuting is at least
	\[
		(1 - q^{-m}) \cdot (1 - q^{-1}) \cdot \Big(1 - \frac{md}{q}\Big) \cdot \frac{1}{2}
		\,\geq\, \Big(1 - \frac{3md}{q}\Big) \cdot \frac{1}{2}\,,
	\]
	and the probability that $\omega$ is commuting is also at least
	$\big(1 - \frac{3md}{q}\big) \cdot \frac{1}{2}$.
\end{lemma}
\begin{proof}
	\uses{def:low-degree-encoding, def:indicator-vector, def:admissible-size, def:binary-representation, lem:downsize_field, lem:one}
	By the definition of the indicator vector
	(Definition~\ref{def:indicator-vector}), expanding the product over the
	coordinates $i \in \{1, \dots, m\}$ into a sum over $y \in \{0,1\}^m$
	gives
	\[
		\mathrm{ind}_m(u_X) \cdot \mathrm{ind}_m(u_Z)
		\,=\, \sum_{y \in \{0,1\}^m} \mathrm{ind}_{m,y}(u_X)\, \mathrm{ind}_{m,y}(u_Z)
		\,=\, \prod_{i=1}^{m} \big( u_{X,i}\, u_{Z,i} + (1 - u_{X,i})(1 - u_{Z,i}) \big)\,,
	\]
	and since $\F_q$ has characteristic $2$, the $i$-th factor equals
	$1 + u_{X,i} + u_{Z,i}$. Hence
	\[
		\gamma(\omega) \,=\, \tr\Big( (r_X r_Z) \cdot \prod_{i=1}^{m} (1 + u_{X,i} + u_{Z,i}) \Big)\,.
	\]
	For any fixed $u_X$, the $m$ factors $1 + u_{X,i} + u_{Z,i}$ are
	independent and uniformly distributed in $\F_q$ over the choice of
	$u_Z$, so all of them are nonzero with probability $(1-q^{-1})^m$;
	independently, $r_X \neq 0$ with probability $1 - q^{-1}$. If $r_X = 0$
	or some factor vanishes, then $\gamma(\omega) = 0$. Otherwise, over the
	uniform choice of $r_Z$, the argument of the trace is a uniformly random
	element of $\F_q$. Since $q = 2^k$ is admissible, the trace of an
	element of $\F_q$ equals the coordinate sum of its binary
	representation: writing $\kappa$ for the binary representation with
	respect to the fixed self-dual normal basis
	(Definition~\ref{def:binary-representation}), item~2 of
	Lemma~\ref{lem:downsize_field} gives
	$\tr(c) = \tr(c \cdot 1) = \kappa(c) \cdot \kappa(1)$, and $\kappa(1)$
	is the all-ones vector by Lemma~\ref{lem:one}. As $\kappa$ is a
	bijection onto $\F_2^k$, a uniformly random element of $\F_q$ has trace
	$0$ with probability exactly $\frac12$. Combining, the probability that
	$\gamma(\omega) \neq 0$ is exactly $(1-q^{-1})^{m+1} \cdot \frac12$, and
	the probability that $\omega$ is commuting is the complement, which is
	at least $\frac12$.

	If $m \leq q$, then
	$(1-q^{-1})^{m+1} \geq (1-q^{-1})^{q+1} = (1-q^{-1})^{q} \cdot (1-q^{-1})$,
	a product of two positive functions of $q$ that are nondecreasing for
	$q \geq 2$, so it is minimized at $q = 2$, where it equals $2^{-3}$;
	hence the anticommuting probability is at least $2^{-4}$.

	Finally, suppose $m, d \geq 1$, and consider the first inequality of
	the display, which bounds the anticommuting probability from below.
	If $md > q$, that lower bound is negative --- it is the product of
	the positive factors $1 - q^{-m}$, $1 - q^{-1}$, and $\frac12$ with
	the negative factor $1 - md/q$ --- so the anticommuting probability,
	being nonnegative, exceeds it. If $md \leq q$, then by Bernoulli's
	inequality
	$(1-q^{-1})^{m} \geq 1 - m/q \geq 1 - md/q \geq (1-q^{-m})(1-md/q)$,
	where the last step multiplies the nonnegative quantity $1 - md/q$
	by $1 - q^{-m} \leq 1$; multiplying by $1 - q^{-1}$ shows that the
	anticommuting probability is at least the first displayed bound. For
	the second inequality of the display, expanding the product shows
	that $(1-q^{-m})(1-q^{-1})(1-md/q)$ exceeds
	$1 - q^{-m} - q^{-1} - md/q$ by
	$q^{-m-1} + \frac{md}{q} \big( q^{-m} + q^{-1} - q^{-m-1} \big)$,
	a sum of nonnegative terms whatever the sign of $1 - md/q$; since
	$md \geq 1$ gives $q^{-m} \leq q^{-1} \leq md/q$, the product is
	therefore at least $1 - 3md/q$.
	The commuting probability is at least
	$\frac12 \geq \big(1 - \frac{3md}{q}\big) \cdot \frac12$.
\end{proof}

\begin{remark}[Source correction: the anticommuting probability]
	\label{rem:omega-anticomm-prob-correction}
	In the source, the statement corresponding to
	Lemma~\ref{fact:omega-anticomm-prob} asserts only the two bounds of the
	final display, without the hypothesis $m, d \geq 1$, and its proof
	begins with the claims that $\mathrm{ind}_m(u_X)$ vanishes if and only
	if $u_X = 0$ and that $u_X = 0$ forces $\omega$ to be commuting. Both
	claims are incorrect: the coordinates of $\mathrm{ind}_m(u)$ sum to $1$
	for every $u \in \F_q^m$, so $\mathrm{ind}_m(u)$ never vanishes, and for
	$(q,m,d) = (2,1,0)$ the tuple $\omega = (0,0,1,1)$ is anticommuting.
	For $d = 0$ the bounds themselves fail: with $(q,m,d) = (2,1,0)$ ---
	a tuple with admissible field size and $m \mid q$, excluded from
	Definition~\ref{def:admissible} only by the requirement $d \geq 1$
	that the source leaves implicit --- the exact anticommuting
	probability
	$(1-q^{-1})^{m+1} \cdot \frac12 = \frac18$ is smaller than the claimed
	lower bound $\big(1 - \frac{3md}{q}\big) \cdot \frac12 = \frac12$. The
	proof given above therefore computes the probability exactly, using the
	characteristic-two product formula for
	$\mathrm{ind}_m(u_X) \cdot \mathrm{ind}_m(u_Z)$, in place of the
	source's argument via the Schwartz--Zippel lemma; the exact value also
	supplies the unconditional constant lower bounds used in the proof of
	Lemma~\ref{lem:qld-win-implications}.
\end{remark}

\begin{definition}[Closeness of question-indexed operators]
	\label{def:approx-question-indexed-operators}
	\uses{def:povm-distance}
	Let $\{A^x\}$ and $\{B^x\}$ be families of operators (not necessarily
	POVM elements) indexed by questions $x \in \mathcal{X}$ but not by
	answers, let $\mu$ be a distribution over $\mathcal{X}$, and let
	$\ket{\psi}$ be a state, both fixed by context. Write
	$A^x \approx_\delta B^x$ to denote
	\[
		\E_{x \sim \mu} \bra{\psi} (A^x - B^x)^\dagger (A^x - B^x) \ket{\psi} \,\leq\, O(\delta)\,.
	\]
	This extends the answer-indexed closeness of
	Definition~\ref{def:povm-distance} to answer-free families of operators,
	such as observables; as there, the right-hand side allows a hidden
	universal constant.
\end{definition}

The following two lemmas relate the closeness between sets of operators and
weighted sums of those operators.

\begin{lemma}
	\label{lem:avg-closeness}
	\uses{def:approx-question-indexed-operators}
	Let $\{A^x\}$ and $\{B^x\}$ be operators indexed by a finite set
	$\mathcal{X}$, let $\mu$ be a probability distribution over
	$\mathcal{X}$, and let $\{\alpha_x\}$ be complex numbers with
	$|\alpha_x| \leq 1$. If $A^x \approx_\delta B^x$ on average over
	$x \sim \mu$, then $A \approx_\delta B$ for the single operators
	$A = \E_{x} \alpha_x A^x$ and $B = \E_{x} \alpha_x B^x$, regarded as
	indexed by a one-element question set; quantitatively,
	% The closeness of Definition def:approx-question-indexed-operators
	% carries a hidden universal constant on both sides; the quantitative
	% sentence records the constant-free inequality the proof establishes.
	\[
		\big\| (A - B) \ket{\psi} \big\|^2
		\,\leq\, \E_x \big\| (A^x - B^x) \ket{\psi} \big\|^2\,,
	\]
	so the conclusion holds with the same implicit constant as the
	hypothesis.
\end{lemma}
\begin{proof}
	Expand
	\[
		\Big\| \E_x \alpha_x (A^x - B^x) \ket{\psi} \Big\|^2
		\,\leq\, \Big( \E_x \big\| (A^x - B^x) \ket{\psi} \big\| \Big)^2
		\,\leq\, \E_x \big\| (A^x - B^x) \ket{\psi} \big\|^2\,,
	\]
	using the triangle inequality followed by Jensen's inequality; the
	right-hand side is at most $O(\delta)$ by the hypothesis, with the same
	implicit constant.
\end{proof}

\begin{lemma}
	\label{lem:povm-to-obs}
	\uses{def:approx-question-indexed-operators, def:povm-distance}
	Let $\mathcal{A}$ be a finite set, let $\{A^x_a\}_{a \in \mathcal{A}}$
	and $\{B^x_a\}_{a \in \mathcal{A}}$ be POVMs, and let
	$\{\alpha_a\}_{a \in \mathcal{A}}$ be complex numbers on the unit
	circle. Let $A^x = \sum_a \alpha_a A^x_a$ and
	$B^x = \sum_a \alpha_a B^x_a$. Then
	\[
		A^x_a \approx_\delta B^x_a \quad\text{implies}\quad A^x \approx_{\delta'} B^x
	\]
	for $\delta' = |\mathcal{A}| \cdot \delta$, where the hypothesis is the
	answer-indexed closeness of Definition~\ref{def:povm-distance} and the
	conclusion is the question-indexed closeness of
	Definition~\ref{def:approx-question-indexed-operators}.
\end{lemma}
\begin{proof}
	\uses{def:povm-distance, def:approx-question-indexed-operators}
	% The source ends this display with the exact bound |A| * delta; under the
	% O(.) reading of def:povm-distance the hypothesis only supplies O(delta),
	% and the constant is retained inside the O(.).
	Expand
	\[
		\E_x \Big\| \sum_a \alpha_a (A^x_a - B^x_a) \ket{\psi} \Big\|^2
		\,\leq\, \E_x \Big( \sum_a \big\| (A^x_a - B^x_a) \ket{\psi} \big\| \Big)^2
		\,\leq\, |\mathcal{A}| \cdot \E_x \sum_a \big\| (A^x_a - B^x_a) \ket{\psi} \big\|^2
		\,\leq\, O\big( |\mathcal{A}| \cdot \delta \big)\,,
	\]
	using the triangle inequality followed by the Cauchy--Schwarz
	inequality; the last step is the hypothesis
	$A^x_a \approx_\delta B^x_a$, whose bound carries a hidden universal
	constant (Definition~\ref{def:povm-distance}), and the conclusion
	holds with the same implicit constant.
\end{proof}

The following lemma shows that POVMs that are approximately projective, in
the sense of being self-consistent, are close to exact projective
measurements.

\begin{lemma}[Orthonormalization lemma]
	\label{lem:ortho}
	\uses{def:consistency, def:povm-distance}
	Let $0 \leq \delta \leq 1$. There exists a function
	$\eta_{\mathrm{ortho}}(\delta) = O\big(\delta^{1/4}\big)$ such that the
	following holds. Let $\ket{\psi}$ be a state on
	$\hilb_{\mathsf{A}} \ot \hilb_{\mathsf{B}}$ where $\hilb_{\mathsf{A}}$
	and $\hilb_{\mathsf{B}}$ are finite dimensional, and let $\{Q_a\}$ and
	$\{R_a\}$ be POVMs on $\hilb_{\mathsf{A}}$ and $\hilb_{\mathsf{B}}$
	respectively such that
	\begin{equation}\label{eq:cons-cond}
		Q_a \ot \Id_{\mathsf{B}} \,\simeq_\delta\, \Id_{\mathsf{A}} \ot R_a\,.
	\end{equation}
	Then there exists a projective measurement $\{P_a\}$ on
	$\hilb_{\mathsf{A}}$ such that, with respect to $\ket{\psi}$,
	\begin{equation}\label{eq:cons-ccl}
		P_a \ot \Id_{\mathsf{B}} \,\approx_{\eta_{\mathrm{ortho}}(\delta)}\, Q_a \ot \Id_{\mathsf{B}}\,.
	\end{equation}
\end{lemma}
\begin{proof}
	This lemma is imported: it first appears in~\cite{KV11}, and a
	self-contained proof is given in the reference [ML20] cited
	in~\cite{Ji2020MIPStar}. The proof is not reproduced here.
\end{proof}

\section{Strategies}
\label{sec:commutation}

% In the source this setup --- fixing the parameter tuple, the strategy and
% its success probability, and making the measurements projective via
% Naimark's theorem --- is untitled prose at the head of the subsection
% (references/qpbt-paper/14_analysis_of_the_pauli_basis_test.tex, lines
% 155-171), not a numbered environment. It is recorded here as a standing
% convention paragraph together with a lemma for the one mathematical step
% it contains, the reduction to a projective strategy.
\paragraph{The strategy under analysis.}
For the remainder of this chapter, as well as in
Chapters~\ref{chap:qpbt-combining} and~\ref{chap:qpbt-extraction}, we fix
an admissible parameter tuple $(q,m,d)$
(Definition~\ref{def:admissible}) and a strategy
$\mathcal{S} = (\ket{\psi}, M)$ for the game
$\mathfrak{G}^{\mathrm{Pauli}}_{(q,m,d)}$ that succeeds with probability
at least $1 - \eps$, for some $\eps \geq 0$, where $\ket{\psi}$ is a
bipartite state on registers $\mathsf{A}$ and $\mathsf{B}$ with associated
Hilbert spaces $\hilb_{\mathsf{A}}$ and $\hilb_{\mathsf{B}}$. The symbol
$M$ denotes the measurement operators of both players; which player is
meant is indicated by the register on which the operator acts (for
example, $M^{(\mathrm{Point},W),u}_a \ot \Id_{\mathsf{B}}$ indicates that
$M^{(\mathrm{Point},W),u}_a$ is viewed as an operator acting on register
$\mathsf{A}$).

The analysis requires the measurements of $\mathcal{S}$ to be projective;
the following lemma shows that this may be assumed without loss of
generality.

\begin{lemma}[Reduction to a projective strategy]
	\label{lem:projective-strategy-setup}
	\uses{def:admissible, def:pauli-question-distribution, def:pauli-win-predicate, def:tensor-product-strategy, def:tensor-product-value, def:projective-strategy-general}
	Let $(q,m,d)$ be an admissible parameter tuple
	(Definition~\ref{def:admissible}) and let
	$\mathcal{S} = (\ket{\psi}, M)$ be a strategy for the game
	$\mathfrak{G}^{\mathrm{Pauli}}_{(q,m,d)}$, where $\ket{\psi}$ is a
	bipartite state on registers $\mathsf{A}$ and $\mathsf{B}$. Then there
	exists a projective strategy
	(Definition~\ref{def:projective-strategy-general}) for
	$\mathfrak{G}^{\mathrm{Pauli}}_{(q,m,d)}$ whose success probability
	equals that of $\mathcal{S}$ and whose state is the extension
	\[
		(\ket{\psi} \ot \ket{0\cdots 0} \ot \ket{0\cdots 0})_{\mathsf{A}\mathsf{B}}
	\]
	of $\ket{\psi}$ by finitely many ancilla qubits initialized in the
	state $\ket{0}$, adjoined to each of the two registers.
\end{lemma}
\begin{proof}
	\uses{thm:naimark, def:tensor-product-value, def:tensor-product-strategy}
	Naimark's theorem (Theorem~\ref{thm:naimark}), applied with the
	families of all of the first player's measurements on register
	$\mathsf{A}$ and all of the second player's on register $\mathsf{B}$,
	yields an auxiliary space and a projective sub-measurement replacing
	each family on each register, together with a product auxiliary state,
	such that every correlation
	$\bra{\psi} M^{x}_a \ot M^{y}_b \ket{\psi}$ of $\mathcal{S}$ is
	reproduced on the extended state. Each auxiliary space embeds into a
	register of finitely many qubits; extend each dilated sub-measurement
	to the enlarged register by adding to one fixed outcome both the
	projector onto the orthogonal complement of the embedded subspace and
	the residual projector of the sub-measurement. The outcome operators
	remain mutually orthogonal projectors, now summing to the identity,
	and no correlation changes: the complement projector annihilates the
	extended state, which is supported in the embedded subspace, and so
	does each residual projector, because for each question pair the
	reproduced correlations sum to one --- the original families being
	POVMs (Definition~\ref{def:tensor-product-strategy}) --- which forces
	the projector $\Id$ minus the sum of each dilated sub-measurement to
	have zero mass on the extended state. Finally, since the auxiliary
	state is a product state, a unitary acting on each player's ancilla
	qubits alone carries it to $\ket{0 \cdots 0} \ot \ket{0 \cdots 0}$;
	conjugating the projective measurements by the same unitaries
	preserves projectivity and all correlations. The success probability
	of a strategy is a linear combination of its correlations, weighted by
	the question distribution and the win predicate
	(Definition~\ref{def:tensor-product-value}), and is therefore
	unchanged.
\end{proof}

In view of Lemma~\ref{lem:projective-strategy-setup} we assume from here
on, replacing $\mathcal{S}$ by the projective strategy provided by that
lemma, that all measurements in $\mathcal{S}$ are projective; the
extended state
$(\ket{\psi} \ot \ket{0\cdots 0} \ot \ket{0\cdots 0})_{\mathsf{A}\mathsf{B}}$
is again denoted $\ket{\psi}$. All statements below refer to this
projective strategy on the extended state, and references to ``the
state $\ket{\psi}$ of the strategy under analysis'' refer to this
convention.

The strategy $\mathcal{S}$ contains measurement operators for each of the
question types and question contents of the Pauli basis test
(Definition~\ref{def:pauli-win-predicate}), such as
$\{M^{(\mathrm{Point},W),u}_a\}_{a \in \F_q}$ for all $u \in \F_q^m$.

\begin{definition}[Strategy observables]
	\label{def:strategy-observables}
	\uses{lem:projective-strategy-setup, def:pauli-win-predicate, def:subfield-trace, def:line, def:line-representative}
	For ``line'' questions we write $\mathrm{Line}$ for either of the
	types $\mathrm{ALine}$ and $\mathrm{DLine}$: a statement about the
	type $\mathrm{Line}$ applies to both types, and the intended type is
	determined by the line description $\ell$. Moreover,
	$\ell$ denotes the description of either an axis-parallel line, given by
	a pair $\ell = (u_0, s)$, or a diagonal line, given by a triple
	$\ell = (u_0, s, v)$ (Definition~\ref{def:line-representative}). For
	every $u \in \F_q^m$, $r \in \F_q$, and $W \in \{X, Z\}$, define
	\begin{equation}\label{eq:qld-strat-obs}
		W^r(u) \,=\, \sum_{a \in \F_q} (-1)^{\tr(ar)} \, M^{(\mathrm{Point},W),u}_a\,.
	\end{equation}
	Since the measurement $\{M^{(\mathrm{Point},W),u}_a\}_{a \in \F_q}$ is
	projective and $(-1)^{\tr(ar)} \in \{\pm 1\}$, the operator $W^r(u)$ is
	an observable with eigenvalues $\pm 1$. It acts on register $\mathsf{A}$
	or $\mathsf{B}$ according to which player's operator
	$M^{(\mathrm{Point},W),u}_a$ denotes.
\end{definition}

The observable $W^r(u)$ is the candidate for the generalized Pauli
observable $\tau^W(r \cdot \mathrm{ind}_m(u))$ that the analysis eventually
extracts from the strategy.

\begin{lemma}[Winning implications]
	\label{lem:qld-win-implications}
	\uses{lem:projective-strategy-setup, def:anticommuting-tuple, def:pauli-question-distribution, def:pauli-win-predicate, def:line-point-dist, def:consistency, def:povm-distance, def:ms-game, def:tensor-product-value, def:low-degree-encoding, def:ideg-deg-polynomials}
	Assume that $\mathcal{S}$ succeeds in $\mathfrak{G}^{\mathrm{Pauli}}_{(q,m,d)}$
	with probability at least $1 - \eps$, for some $\eps \geq 0$. Then the
	following hold, with respect to the state $\ket{\psi}$ of the strategy
	under analysis (Lemma~\ref{lem:projective-strategy-setup} and the
	convention following it).
	\begin{enumerate}
		\item \label{enu:qld-cons} (\textbf{Consistency check}) On average
		over $(\mathsf{t},x)$ sampled from the question distribution of
		$\mathfrak{G}^{\mathrm{Pauli}}_{(q,m,d)}$
		(Definition~\ref{def:pauli-question-distribution}), we have
		$M^{\mathsf{t},x}_a \ot \Id_{\mathsf{B}} \simeq_\eps \Id_{\mathsf{A}} \ot M^{\mathsf{t},x}_a$.
		\item \label{enu:qld-ld} (\textbf{Low-degree check}) For all
		$W \in \{X,Z\}$, on average over $(\ell,u)$ drawn from the
		line-point distribution (Definition~\ref{def:line-point-dist}), we
		have
		% The source states this item with answer summation over F_q only,
		% implicitly reading the evaluation classes of the line measurement
		% as a complete measurement; see the paragraph after the item list
		% and the remark rem:qld-win-implications-typos.
		$M^{(\mathrm{Line},W),\ell}_{[\mathrm{eval}_u(\cdot)=a]} \ot \Id_{\mathsf{B}} \simeq_\eps \Id_{\mathsf{A}} \ot M^{(\mathrm{Point},W),u}_a$,
		where the answer summation is over the completed outcome set
		$a \in \F_q \cup \{\bot\}$ described after the item list.
		\item \label{enu:qld-pauli-basis-cons} (\textbf{Pauli basis
		consistency check}) For all $W \in \{X,Z\}$, on average over
		uniformly random $u \in \F_q^m$, we have
		$M^{(\mathrm{Point},W),u}_a \ot \Id_{\mathsf{B}} \simeq_\eps \Id_{\mathsf{A}} \ot M^{(\mathrm{Pauli},W)}_{[g_h(u)=a]}$,
		where
		$M^{(\mathrm{Pauli},W)}_{[g_h(u)=a]} = \sum_{h \in \F_q^{M} \colon g_h(u)=a} M^{(\mathrm{Pauli},W)}_h$
		and $g_h$ is the low-degree encoding of $h$
		(see Definition~\ref{def:low-degree-encoding}).
		\item \label{enu:qld-comm} (\textbf{Commutation check}) For all
		$W \in \{X,Z\}$, on average over commuting
		$\omega = (u_X,u_Z,r_X,r_Z)$ (that is, over a uniformly random
		$\omega$ conditioned on being commuting), we have
		$M^{(\mathrm{Pair},W),\omega}_a \ot \Id_{\mathsf{B}} \simeq_\eps \Id_{\mathsf{A}} \ot M^{\mathrm{Pair},\omega}_{[\beta_W=a]}$,
		where the answer summation is over $a \in \F_2$.
		\item \label{enu:qld-comm-cons} (\textbf{Commutation consistency
		check}) For all $W \in \{X,Z\}$, on average over commuting
		$\omega$, we have
		$M^{(\mathrm{Point},W),u_W}_{[\tr(\cdot\,r_W)=a]} \ot \Id_{\mathsf{B}} \simeq_\eps \Id_{\mathsf{A}} \ot M^{(\mathrm{Pair},W),\omega}_a$,
		where the answer summation is over $a \in \F_2$.
		\item \label{enu:qld-ms} (\textbf{Magic square check}) For any
		anticommuting $\omega$, let $\Lambda_\omega$ be the value of the
		strategy
		$\big( \ket{\psi}, \{M^{\mathrm{Constraint}_i,\omega}_{\alpha_1,\alpha_2,\alpha_3}\} \cup \{M^{\mathrm{Variable}_j,\omega}_a\} \big)$
		in the Magic Square game $\mathfrak{G}^{\mathrm{MS}}$
		(Definition~\ref{def:ms-game}). Then
		$\E_\omega \Lambda_\omega = 1 - O(\eps)$, where the expectation is
		over a uniformly random $\omega$ conditioned on being
		anticommuting.
		\item \label{enu:qld-ms-cons} (\textbf{Magic square consistency
		check}) For all $W \in \{X,Z\}$, on average over anticommuting
		$\omega$,
		\begin{gather*}
			M^{(\mathrm{Point},X),u_X}_{[\tr(\cdot\,r_X)=a]} \ot \Id_{\mathsf{B}} \,\simeq_\eps\, \Id_{\mathsf{A}} \ot M^{\mathrm{Variable}_1,\omega}_a\,,\\
			M^{(\mathrm{Point},Z),u_Z}_{[\tr(\cdot\,r_Z)=a]} \ot \Id_{\mathsf{B}} \,\simeq_\eps\, \Id_{\mathsf{A}} \ot M^{\mathrm{Variable}_5,\omega}_a\,,
		\end{gather*}
		where the answer summations are over $\F_2$.
	\end{enumerate}
	Here
	$M^{(\mathrm{Point},W),u}_{[\tr(\cdot\,r)=a]} = \sum_{a' \in \F_q \colon \tr(a'r)=a} M^{(\mathrm{Point},W),u}_{a'}$
	for $a \in \F_2$; in item~\ref{enu:qld-ld} the two families are the
	following POVMs indexed by $\F_q \cup \{\bot\}$, as the relation
	$\simeq_\eps$ of Definition~\ref{def:consistency} requires: on the left,
	the completed family of evaluation classes of
	Definition~\ref{def:ideg-deg-polynomials}, whose class
	$M^{(\mathrm{Line},W),\ell}_{[\mathrm{eval}_u(\cdot)=\bot]}$ carries the
	line answers with no evaluation at $u$ and is nonzero only for a
	zero-direction line, where the classes indexed by $a \in \F_q$ form
	only a sub-measurement; on the right, the point measurement completed
	by the zero operator $M^{(\mathrm{Point},W),u}_{\bot} = 0$; the answer
	to a question of type $\mathrm{Pair}$ is a
	pair $(\beta_X,\beta_Z) \in \F_2 \times \F_2$, and
	$M^{\mathrm{Pair},\omega}_{[\beta_W=a]}$ denotes its coarse-graining
	onto the $W$-component. Furthermore, all of the approximations above
	also hold with $\approx_\eps$ in place of $\simeq_\eps$, and with the
	tensor factors interchanged.
\end{lemma}
\begin{proof}
	\uses{fact:omega-anticomm-prob, fact:agreement, def:admissible, def:pauli-question-distribution, def:pauli-win-predicate, def:ld-win-predicate, def:ideg-deg-polynomials, def:consistency}
	The question types are sampled according to the type graph of the Pauli
	basis test (Definition~\ref{def:pauli-question-distribution}), which has
	constant size, so each subtest of the win predicate
	(Definition~\ref{def:pauli-win-predicate}) is executed with probability
	bounded below by a universal constant.
	Items~\ref{enu:qld-cons},~\ref{enu:qld-ld}
	and~\ref{enu:qld-pauli-basis-cons} then follow directly from the
	corresponding subtests and the assumption that $\mathcal{S}$ succeeds
	with probability at least $1-\eps$, where item~\ref{enu:qld-ld} follows
	from the consistency subtest of the low individual degree game executed
	inside the Pauli basis test. In more detail for item~\ref{enu:qld-ld}:
	the disagreement mass of the two completed families is rejection mass
	of the line-versus-point clauses
	(Definition~\ref{def:ld-win-predicate}, invoked by the low-degree
	check of Definition~\ref{def:pauli-win-predicate}). An outcome pair
	$(a,b)$ with $a \neq b$ in $\F_q$ is a line answer evaluating at $u$
	to $a$ against the point answer $b$, which the clause rejects; a pair
	$(\bot,b)$ is a line answer with no evaluation at $u$, which the
	universally quantified clause rejects against every point answer $b$
	(Remark~\ref{rem:ld-win-zero-direction} and
	Remark~\ref{rem:deg-line-representatives}); and the pairs $(a,\bot)$
	contribute no mass, since the class
	$M^{(\mathrm{Point},W),u}_{\bot}$ is the zero operator. The subtests underlying
	items~\ref{enu:qld-comm} and~\ref{enu:qld-comm-cons} accept vacuously
	unless the sampled tuple $\omega$ is commuting, and those underlying
	items~\ref{enu:qld-ms} and~\ref{enu:qld-ms-cons} accept vacuously
	unless $\omega$ is anticommuting. Since the parameter tuple $(q,m,d)$
	is admissible (Definition~\ref{def:admissible}), $m$ divides $q$, so
	$1 \leq m \leq q$, and Lemma~\ref{fact:omega-anticomm-prob} shows that
	a uniformly sampled tuple $\omega$ is commuting with probability at
	least $\frac12$ and anticommuting with probability at least $2^{-4}$.
	Conditioning on either event therefore increases the failure
	probability of the corresponding subtest by a factor of at most $2^4$,
	and items~\ref{enu:qld-comm},~\ref{enu:qld-comm-cons}
	and~\ref{enu:qld-ms-cons} follow from the corresponding subtests.
	Item~\ref{enu:qld-ms} follows in the same way, because the Magic Square
	check is executed with constant probability over the choice of a pair
	of questions and the anticommuting event has probability at least
	$2^{-4}$. Finally, that all approximations
	also hold with $\approx_\eps$ in place of $\simeq_\eps$ is immediate
	from Lemma~\ref{fact:agreement}.
\end{proof}

\begin{lemma}[Winning implications for observables]
	\label{lem:qld-win-implications-obs}
	% The source labels this lemma `len:qld-win-implications-obs'; the prefix
	% `len' is a typo for `lem'. The label above is the corrected form.
	\uses{def:strategy-observables, def:approx-question-indexed-operators, def:anticommuting-tuple, lem:projective-strategy-setup}
	The following hold for the observables $W^r(u)$ of
	Definition~\ref{def:strategy-observables}, with respect to the state
	$\ket{\psi}$ of the strategy under analysis
	(Lemma~\ref{lem:projective-strategy-setup} and the convention
	following it).
	\begin{enumerate}
		\item For all $W \in \{X,Z\}$ and $r \in \F_q$, on average over
		uniformly random $u \in \F_q^m$,
		\begin{equation}\label{eq:pts-obs-consistency}
			W^r(u) \ot \Id_{\mathsf{B}} \,\approx_\eps\, \Id_{\mathsf{A}} \ot W^r(u)\,.
		\end{equation}
		\item On average over uniformly random
		$\omega = (u_X,u_Z,r_X,r_Z) \in (\F_q^m)^2 \times (\F_q)^2$,
		\begin{equation}\label{eq:pts-obs-commutation}
			X^{r_X}(u_X)\, Z^{r_Z}(u_Z) \ot \Id_{\mathsf{B}}
			\,\approx_{\sqrt{\eps}}\,
			(-1)^{\gamma(\omega)}\, Z^{r_Z}(u_Z)\, X^{r_X}(u_X) \ot \Id_{\mathsf{B}}\,,
		\end{equation}
		where
		$\gamma(\omega) = \tr((r_X \cdot \mathrm{ind}_m(u_X)) \cdot (r_Z \cdot \mathrm{ind}_m(u_Z)))$
		as in Definition~\ref{def:anticommuting-tuple}.
	\end{enumerate}
	Furthermore, all of the approximations above also hold with the tensor
	factors interchanged.
\end{lemma}
\begin{proof}
	\uses{lem:qld-win-implications, lem:povm-to-obs, fact:data-processing, fact:agreement, lem:commutation-analysis, thm:ms-rigidity, fact:add-a-proj, def:strategy-observables}
	For~\eqref{eq:pts-obs-consistency}: item~\ref{enu:qld-cons} of
	Lemma~\ref{lem:qld-win-implications} and the data processing inequality
	(Lemma~\ref{fact:data-processing}), using that the type
	$(\mathrm{Point},W)$ is sampled with constant probability, give the
	self-consistency of the coarse-grained measurements
	$\{M^{(\mathrm{Point},W),u}_{[\tr(\cdot\,r)=a]}\}_{a \in \F_2}$ on
	average over $u$; Lemma~\ref{lem:povm-to-obs}, applied with
	$\mathcal{A} = \F_2$ and $\alpha_a = (-1)^a$, converts this
	to~\eqref{eq:pts-obs-consistency}. For the commuting case
	of~\eqref{eq:pts-obs-commutation}: chaining
	items~\ref{enu:qld-comm-cons},~\ref{enu:qld-comm}
	and~\ref{enu:qld-cons} of Lemma~\ref{lem:qld-win-implications} via
	Lemma~\ref{fact:agreement} gives, on average over commuting $\omega$,
	$M^{(\mathrm{Point},W),u_W}_{[\tr(\cdot\,r_W)=a]} \ot \Id_{\mathsf{B}} \approx_\eps \Id_{\mathsf{A}} \ot M^{\mathrm{Pair},\omega}_{[\beta_W=a]}$;
	since the measurement
	$\{M^{\mathrm{Pair},\omega}_{\beta_X,\beta_Z}\}$ is projective,
	Lemma~\ref{lem:commutation-analysis}, applied with the $X$- and
	$Z$-point projections and the $\mathrm{Pair}$ measurement, yields the
	approximate commutation of the projections
	$M^{(\mathrm{Point},X),u_X}_{[\tr(\cdot\,r_X)=b]}$ and
	$M^{(\mathrm{Point},Z),u_Z}_{[\tr(\cdot\,r_Z)=c]}$; expanding
	$W^r(u) = \Id - 2 M^{(\mathrm{Point},W),u}_{[\tr(\cdot\,r)=1]}$
	transfers this to the observables with error $\eps$. For the
	anticommuting case: item~\ref{enu:qld-ms-cons} of
	Lemma~\ref{lem:qld-win-implications} together with
	Lemma~\ref{lem:povm-to-obs} relates $X^{r_X}(u_X)$ and $Z^{r_Z}(u_Z)$
	to the Magic Square observables
	$M^{\mathrm{Variable}_j,\omega} = M^{\mathrm{Variable}_j,\omega}_0 - M^{\mathrm{Variable}_j,\omega}_1$
	for $j \in \{1,5\}$; for each anticommuting $\omega$, the Magic Square
	rigidity theorem (Theorem~\ref{thm:ms-rigidity}) applied with
	$\eps_\omega = 1 - \Lambda_\omega$ gives
	$\Id \ot M^{\mathrm{Variable}_1,\omega} M^{\mathrm{Variable}_5,\omega} \approx_{\sqrt{\eps_\omega}} - \Id \ot M^{\mathrm{Variable}_5,\omega} M^{\mathrm{Variable}_1,\omega}$;
	averaging over anticommuting $\omega$ with Jensen's inequality
	($\E \eps_\omega^{1/2} \leq (\E \eps_\omega)^{1/2} = O(\eps^{1/2})$,
	using item~\ref{enu:qld-ms}) and transferring across the tensor factor
	with Lemma~\ref{fact:add-a-proj}, applied twice, yields the
	anticommutation with error $\sqrt{\eps}$. The two cases combine into
	the $(-1)^{\gamma(\omega)}$-signed statement over uniformly random
	$\omega$, since the uniform average is bounded by the sum of the two
	conditional averages.
\end{proof}

\begin{remark}[Source corrections]
	\label{rem:qld-win-implications-typos}
	Two typographical slips in the source are corrected above. In
	item~\ref{enu:qld-ms} of Lemma~\ref{lem:qld-win-implications} the
	source writes the expectation over $\omega$ subject to the condition
	$\tr((r_X \cdot u_X) \cdot (r_Z \cdot u_Z)) \neq 0$, omitting the map
	$\mathrm{ind}_m$; the intended condition, consistent with the
	surrounding text and with Definition~\ref{def:anticommuting-tuple}, is
	$\gamma(\omega) \neq 0$, that is, $\omega$ anticommuting, and this is
	the form stated above. In the source proof of
	Lemma~\ref{lem:qld-win-implications-obs}, one display reads
	$Z^{r_X}(u_X)\,X^{r_Z}(u_Z)$ where $Z^{r_Z}(u_Z)\,X^{r_X}(u_X)$ is
	meant.
	% A third slip in the source, the label prefix `len' in place of `lem'
	% on the lemma corresponding to lem:qld-win-implications-obs, is
	% recorded in the comment at that lemma; no alias label is defined.
	In addition, the source proof of Lemma~\ref{lem:qld-win-implications}
	restricts attention to the case $6md \leq q$ and disposes of the
	remaining parameters by an appeal to Theorem~\ref{thm:pauli}, the
	theorem whose proof the lemma is part of; that argument does not
	establish the unrestricted statement of the lemma, and for $d = 0$ the
	probability bound it invokes is false
	(Remark~\ref{rem:omega-anticomm-prob-correction}). The proof given
	above instead uses the lower bounds of
	Lemma~\ref{fact:omega-anticomm-prob} on the probabilities of the
	commuting and anticommuting events, which are universal constants for
	every admissible parameter tuple, so no restriction on the parameters
	is needed.
	Finally, the source states item~\ref{enu:qld-ld} of
	Lemma~\ref{lem:qld-win-implications} with answer summation over
	$a \in \F_q$ only, implicitly reading the evaluation classes of the
	line measurement as a complete measurement. Under the coefficient
	reading of Definition~\ref{def:ideg-deg-polynomials}
	(Remark~\ref{rem:deg-line-representatives}) the classes indexed by
	$a \in \F_q$ form only a sub-measurement when the line has zero
	direction --- a case the question distribution produces with positive
	probability (Remark~\ref{rem:ld-win-zero-direction}) --- while the
	relation $\simeq_\eps$ of Definition~\ref{def:consistency} is defined
	for POVMs. Item~\ref{enu:qld-ld} therefore completes both families by
	an explicit $\bot$ outcome; the same implicit completeness assumption
	occurs in the source proof of the statement corresponding to
	Lemma~\ref{lem:qld-comm-line-cons}, whose proof below invokes the
	completed families instead.
\end{remark}

\section{Expanding the Hilbert space and defining commuting observables}
\label{sec:expanding}

The next step of the analysis adjoins ancillary EPR registers to the
players' state and lifts the strategy observables and measurements to the
expanded state. The point of the construction is that the ancilla factor of
each lifted observable carries the \emph{exact} Pauli commutation sign, so
that the signs cancel and the lifted observables approximately commute for
every tuple $\omega$.

\begin{definition}[Expanded state]
	\label{def:expanded-state}
	\uses{lem:projective-strategy-setup, def:EPR}
	Recall that $M = 2^m$. Introduce registers $\mathsf{A}', \mathsf{A}'',
	\mathsf{B}', \mathsf{B}''$, each with Hilbert space isomorphic to
	$(\C^q)^{\ot M}$, and define the state
	\begin{equation}\label{eq:def-psihat}
		\ket{\hat{\psi}}_{\mathsf{A}\mathsf{A}'\mathsf{A}''\mathsf{B}\mathsf{B}'\mathsf{B}''}
		\,=\, (\ket{\psi})_{\mathsf{A}\mathsf{B}}
		\ot \ket{\mathrm{EPR}_q}^{\ot M}_{\mathsf{A}'\mathsf{A}''}
		\ot \ket{\mathrm{EPR}_q}^{\ot M}_{\mathsf{B}'\mathsf{B}''}\,,
	\end{equation}
	where $\ket{\psi}$ is the extended strategy state of the strategy
	under analysis (Lemma~\ref{lem:projective-strategy-setup} and the
	convention following it) and
	$\ket{\mathrm{EPR}_q}$ is the maximally entangled state on
	$\C^q \ot \C^q$ (Definition~\ref{def:EPR}). For the remainder of this
	chapter, as well as in Chapters~\ref{chap:qpbt-combining}
	and~\ref{chap:qpbt-extraction}, all approximations $\approx$ and
	$\simeq$ are evaluated with respect to the state $\ket{\hat{\psi}}$,
	unless stated otherwise.
\end{definition}

\begin{definition}[Expanded observables]
	\label{def:expanded-observables}
	\uses{def:strategy-observables, def:expanded-state, lem:twisted-commutation}
	For all $W \in \{X,Z\}$, $r \in \F_q$, and $u \in \F_q^m$, define the
	observable
	\begin{equation}\label{eq:lc-23}
		\hat{W}^r(u) \,=\, W^r(u) \ot \tau^W(r \cdot \mathrm{ind}_m(u))\,,
	\end{equation}
	where $W^r(u)$ is the strategy observable of~\eqref{eq:qld-strat-obs}
	and $\tau^W(r \cdot \mathrm{ind}_m(u))$ is the generalized Pauli
	observable acting on $(\C^q)^{\ot M}$, introduced in
	Chapter~\ref{chap:qpbt-algebra}
	(see Lemma~\ref{lem:twisted-commutation}). As a product
	of two observables with eigenvalues $\pm 1$ acting on disjoint
	registers, $\hat{W}^r(u)$ is itself an observable with eigenvalues
	$\pm 1$. If $W^r(u)$ is viewed as acting on register $\mathsf{A}$, then
	$\hat{W}^r(u)$ is viewed as acting on registers
	$\mathsf{A}\mathsf{A}'$, with the $\tau^W$ factor acting on
	$\mathsf{A}'$. The registers are generally indicated by subscripts, as
	in $(\hat{W}^r(u))_{\mathsf{A}\mathsf{A}'}$ or
	$(\hat{W}^r(u))_{\mathsf{B}\mathsf{A}''}$, the latter having the
	strategy factor on $\mathsf{B}$ and the Pauli factor on
	$\mathsf{A}''$.
\end{definition}

\begin{definition}[Expanded point measurements]
	\label{def:expanded-point-measurement}
	\uses{def:expanded-observables, lem:pauli-observable-expansion, def:low-degree-encoding, def:subfield-trace}
	For $W \in \{X,Z\}$, $u \in \F_q^m$ and $a \in \F_q$, define
	\[
		\hat{M}^{(\mathrm{Point},W),u}_a \,=\, \E_{r \in \F_q} (-1)^{\tr(ar)} \, \hat{W}^r(u)\,,
	\]
	and define the projector
	\begin{equation}\label{eq:qld-point-obs-def}
		\tau^{W,u}_a \,=\, \tau^W_{[g_\cdot(u) = a]}
		\,=\, \sum_{h \in \F_q^{M} \colon h \cdot \mathrm{ind}_m(u) = a} \tau^W_h\,,
	\end{equation}
	where $\{\tau^W_h\}_{h \in \F_q^{M}}$ are the joint eigenprojectors of
	the generalized Pauli $W$-observables
	(see Lemma~\ref{lem:pauli-observable-expansion}).
	Expanding both tensor factors of $\hat{W}^r(u)$
	by~\eqref{eq:qld-strat-obs} and by the Fourier expansion of the
	generalized Pauli observables
	(see Lemma~\ref{lem:pauli-observable-expansion}), and
	averaging over $r$, gives the equivalent convolution form
	\[
		\hat{M}^{(\mathrm{Point},W),u}_a
		\,=\, \sum_{\substack{a',a'' \in \F_q \\ a' + a'' = a}}
		M^{(\mathrm{Point},W),u}_{a'} \ot \tau^{W,u}_{a''}\,.
	\]
	In particular each $\hat{M}^{(\mathrm{Point},W),u}_a$ is a projection
	and $\{\hat{M}^{(\mathrm{Point},W),u}_a\}_{a \in \F_q}$ is a projective
	measurement.
\end{definition}

\begin{definition}[Trace-coarse-grained point projections]
	\label{def:expanded-point-trace-projection}
	\uses{def:expanded-point-measurement, def:expanded-observables, def:subfield-trace}
	For fixed $u \in \F_q^m$, $r \in \F_q$ and $b \in \F_2$, define the
	projection
	\begin{equation}\label{eq:qld-def-mptur}
		\hat{M}^{(\mathrm{Point},W),u,r}_b \,=\, \sum_{a \in \F_q \colon \tr(ar) = b} \hat{M}^{(\mathrm{Point},W),u}_a\,,
	\end{equation}
	which can be equivalently written as
	\begin{equation}\label{eq:lc-22}
		\hat{M}^{(\mathrm{Point},W),u,r}_b \,=\, \frac{1}{2} \Big( \Id + (-1)^b\, \hat{W}^r(u) \Big)\,.
	\end{equation}
	The register conventions are those of
	Definition~\ref{def:expanded-observables}.
\end{definition}

\begin{definition}[Symmetric equivalents]
	\label{def:symmetric-equivalents}
	\uses{def:expanded-state}
	For the remainder of this chapter, as well as in
	Chapters~\ref{chap:qpbt-combining} and~\ref{chap:qpbt-extraction}, the
	six registers
	$\mathsf{A}\mathsf{A}'\mathsf{A}''\mathsf{B}\mathsf{B}'\mathsf{B}''$
	are partitioned into two ``parties'' in two different ways. The first
	scheme takes $\mathsf{A}\mathsf{A}'$ as one party and
	$\mathsf{B}\mathsf{A}''$ as the other; the second scheme takes
	$\mathsf{B}\mathsf{B}'$ as the first party and
	$\mathsf{A}\mathsf{B}''$ as the second. The \emph{symmetric
	equivalents} of a bipartite relation of the form
	$M_{\mathsf{A}\mathsf{A}'} \approx_\delta N_{\mathsf{B}\mathsf{A}''}$
	are the three relations
	\[
		N_{\mathsf{A}\mathsf{A}'} \approx_\delta M_{\mathsf{B}\mathsf{A}''}\,,\qquad
		M_{\mathsf{B}\mathsf{B}'} \approx_\delta N_{\mathsf{A}\mathsf{B}''}\,,\qquad
		N_{\mathsf{B}\mathsf{B}'} \approx_\delta M_{\mathsf{A}\mathsf{B}''}\,,
	\]
	obtained by passing to the second bipartitioning scheme and by
	interchanging the first and second party; the symmetric equivalents of
	a single-party relation of the form
	$M_{\mathsf{A}\mathsf{A}'} \approx N_{\mathsf{A}\mathsf{A}'}$ are the
	same relation on the register pairs $\mathsf{B}\mathsf{A}''$,
	$\mathsf{B}\mathsf{B}'$ and $\mathsf{A}\mathsf{B}''$.
\end{definition}

\begin{lemma}[Transfer to symmetric equivalents]
	\label{lem:symmetric-equivalents-transfer}
	\notready
	\uses{def:symmetric-equivalents, def:expanded-state, lem:projective-strategy-setup, def:symmetric-game, def:EPR, def:povm-distance}
	Let $\mathcal{S} = (\ket{\psi}, M)$ be the projective strategy under
	analysis (Lemma~\ref{lem:projective-strategy-setup} and the convention
	following it), assumed symmetric in the sense of
	Definition~\ref{def:symmetric-game}: the state $\ket{\psi}$ is
	invariant under exchanging registers $\mathsf{A}$ and $\mathsf{B}$, and
	both players measure the same operators. Let $\ket{\hat{\psi}}$ be the
	expanded state~\eqref{eq:def-psihat}, let
	\[
		\sigma \colon\;
		\mathsf{A} \leftrightarrow \mathsf{B},\quad
		\mathsf{A}' \leftrightarrow \mathsf{B}',\quad
		\mathsf{A}'' \leftrightarrow \mathsf{B}''\,,
		\qquad\qquad
		\theta \colon\;
		\mathsf{A} \leftrightarrow \mathsf{B},\quad
		\mathsf{A}' \leftrightarrow \mathsf{A}''\,,
	\]
	be the two permutations of the six registers
	$\mathsf{A}\mathsf{A}'\mathsf{A}''\mathsf{B}\mathsf{B}'\mathsf{B}''$
	given by these transpositions, and let $U_\sigma$ and $U_\theta$ be the
	corresponding permutation unitaries. Write
	$\Pi = \{\Id,\, U_\sigma,\, U_\theta,\, U_\sigma U_\theta\}$ and let
	\[
		\mathcal{R} = \{\mathsf{A}\mathsf{A}',\; \mathsf{B}\mathsf{A}'',\;
		\mathsf{B}\mathsf{B}',\; \mathsf{A}\mathsf{B}''\}
	\]
	be the four register placements of
	Definition~\ref{def:symmetric-equivalents}. Then:
	\begin{enumerate}
		\item Every $U \in \Pi$ fixes $\ket{\hat{\psi}}$, and the induced
		action of $\Pi$ on $\mathcal{R}$ is simply transitive.
		\item Let $P$ and $Q$ be operators on
		$\hilb \ot (\C^q)^{\ot M}$, where $\hilb$ is the common Hilbert
		space of registers $\mathsf{A}$ and $\mathsf{B}$, let
		$\mathsf{C}\mathsf{C}', \mathsf{D}\mathsf{D}' \in \mathcal{R}$ and
		let $U \in \Pi$. Then
		\[
			\bigl\| \bigl( (P)_{\mathsf{C}\mathsf{C}'} - (Q)_{\mathsf{D}\mathsf{D}'} \bigr) \ket{\hat{\psi}} \bigr\|
			\,=\,
			\bigl\| \bigl( (P)_{U(\mathsf{C}\mathsf{C}')} - (Q)_{U(\mathsf{D}\mathsf{D}')} \bigr) \ket{\hat{\psi}} \bigr\|\,,
		\]
		the placement notation being that of
		Definition~\ref{def:expanded-observables}. Consequently, for
		families $\{P^x_a\}$ and $\{Q^x_a\}$ indexed by a question $x$ drawn
		from a distribution and an answer $a$, and for every $\delta \ge 0$,
		$(P^x_a)_{\mathsf{C}\mathsf{C}'} \approx_\delta (Q^x_a)_{\mathsf{D}\mathsf{D}'}$
		holds on $\ket{\hat{\psi}}$ if and only if
		$(P^x_a)_{U(\mathsf{C}\mathsf{C}')} \approx_\delta (Q^x_a)_{U(\mathsf{D}\mathsf{D}')}$
		does, with the same $\delta$
		(Definition~\ref{def:povm-distance}).
	\end{enumerate}
	In particular, taking $\mathsf{C}\mathsf{C}' = \mathsf{A}\mathsf{A}'$
	and $\mathsf{D}\mathsf{D}' = \mathsf{B}\mathsf{A}''$, respectively
	$\mathsf{D}\mathsf{D}' = \mathsf{A}\mathsf{A}'$, the images under
	$U_\theta$, $U_\sigma$ and $U_\sigma U_\theta$ are exactly the three
	symmetric equivalents of Definition~\ref{def:symmetric-equivalents},
	with the same error parameter.
\end{lemma}
\begin{proof}
	\uses{def:expanded-state, def:EPR, def:symmetric-game, def:povm-distance}
	Each of $\sigma$ and $\theta$ transposes registers carrying isomorphic
	Hilbert spaces --- $\mathsf{A}$ and $\mathsf{B}$ both carry $\hilb$,
	and the four ancilla registers are all isomorphic to
	$(\C^q)^{\ot M}$ (Definition~\ref{def:expanded-state}) --- so the
	permutation unitaries $U_\sigma$ and $U_\theta$ are defined and satisfy
	$U (P)_{\mathsf{C}\mathsf{C}'} U^\dagger = (P)_{U(\mathsf{C})U(\mathsf{C}')}$
	for every operator $P$ and every placement.

	For item~1, both generators fix $\ket{\hat{\psi}} = (\ket{\psi})_{\mathsf{A}\mathsf{B}} \ot \ket{\mathrm{EPR}_q}^{\ot M}_{\mathsf{A}'\mathsf{A}''} \ot \ket{\mathrm{EPR}_q}^{\ot M}_{\mathsf{B}'\mathsf{B}''}$:
	the factor $(\ket{\psi})_{\mathsf{A}\mathsf{B}}$ is invariant under
	$\mathsf{A} \leftrightarrow \mathsf{B}$ by the symmetry hypothesis;
	$\sigma$ exchanges the two ancilla factors, which are identical; and
	$\theta$ leaves $\ket{\mathrm{EPR}_q}^{\ot M}_{\mathsf{B}'\mathsf{B}''}$
	untouched and exchanges the two halves of each pair in
	$\ket{\mathrm{EPR}_q}^{\ot M}_{\mathsf{A}'\mathsf{A}''}$, under which
	$\ket{\mathrm{EPR}_q}$ is invariant (Definition~\ref{def:EPR}). The
	action on $\mathcal{R}$ sends $\mathsf{A}\mathsf{A}'$ to
	$\mathsf{A}\mathsf{A}'$, $\mathsf{B}\mathsf{B}'$,
	$\mathsf{B}\mathsf{A}''$ and $\mathsf{A}\mathsf{B}''$ under $\Id$,
	$U_\sigma$, $U_\theta$ and $U_\sigma U_\theta$ respectively, so it is
	simply transitive on the four placements.

	For item~2, insert $U^\dagger U$ and use $U \ket{\hat{\psi}} = \ket{\hat{\psi}}$
	together with unitary invariance of the norm:
	\[
		\bigl\| \bigl( (P)_{\mathsf{C}\mathsf{C}'} - (Q)_{\mathsf{D}\mathsf{D}'} \bigr) \ket{\hat{\psi}} \bigr\|
		= \bigl\| U \bigl( (P)_{\mathsf{C}\mathsf{C}'} - (Q)_{\mathsf{D}\mathsf{D}'} \bigr) U^\dagger \, U \ket{\hat{\psi}} \bigr\|
		= \bigl\| \bigl( (P)_{U(\mathsf{C}\mathsf{C}')} - (Q)_{U(\mathsf{D}\mathsf{D}')} \bigr) \ket{\hat{\psi}} \bigr\|\,.
	\]
	Summing the squares over the answers $a$ and averaging over the
	questions $x$ preserves the identity, so the two sides of the asserted
	equivalence of $\approx_\delta$ relations
	(Definition~\ref{def:povm-distance}) are equal quantities. The final
	assertion is item~1 applied to the placements listed there.
\end{proof}

Several lemmas below conclude by asserting that all symmetric equivalents
of the stated approximations also hold, in each case by
Lemma~\ref{lem:symmetric-equivalents-transfer}; later chapters invoke
exactly these variants.

\begin{lemma}[Consistency and approximate commutation of the expanded point measurements]
	\label{lem:qld-comm-cons}
	\uses{def:expanded-point-measurement, def:expanded-point-trace-projection, def:expanded-state, def:approx-question-indexed-operators, def:povm-distance, def:symmetric-equivalents}
	There exists a function $\delta_{\mathrm{Anticom}}(\eps) = \poly(\eps)$
	such that the following holds, with respect to the state
	$\ket{\hat{\psi}}$ of Definition~\ref{def:expanded-state}.
	\begin{enumerate}
		\item (\textbf{Self-consistency}) For all $W \in \{X,Z\}$, on
		average over $r \in \F_q$ and $u \in \F_q^m$ (the relation does not
		depend on $r$), we have
		\[
			(\hat{M}^{(\mathrm{Point},W),u}_a)_{\mathsf{A}\mathsf{A}'}
			\,\approx_\eps\,
			(\hat{M}^{(\mathrm{Point},W),u}_a)_{\mathsf{B}\mathsf{A}''}\,,
		\]
		where the answer summation is over $a \in \F_q$.
		\item (\textbf{Approximate commutation}) For all $b, b' \in \F_2$,
		on average over uniformly random
		$\omega = (u_X,u_Z,r_X,r_Z) \in (\F_q^m)^2 \times (\F_q)^2$, we
		have
		\[
			(\hat{M}^{(\mathrm{Point},X),u_X,r_X}_{b}\,
			\hat{M}^{(\mathrm{Point},Z),u_Z,r_Z}_{b'})_{\mathsf{A}\mathsf{A}'}
			\,\approx_{\delta_{\mathrm{Anticom}}}\,
			(\hat{M}^{(\mathrm{Point},Z),u_Z,r_Z}_{b'}\,
			\hat{M}^{(\mathrm{Point},X),u_X,r_X}_{b})_{\mathsf{A}\mathsf{A}'}\,.
		\]
	\end{enumerate}
	Furthermore, all symmetric equivalents of these approximations also
	hold, and the proof yields $\delta_{\mathrm{Anticom}}(\eps) = \sqrt{\eps}$.
\end{lemma}
\begin{proof}
	\uses{lem:qld-win-implications, lem:qld-win-implications-obs, fact:data-processing, lem:twisted-commutation, def:expanded-point-measurement, def:expanded-point-trace-projection, def:expanded-observables, lem:symmetric-equivalents-transfer}
	For item~1: the original point measurements
	$\{M^{(\mathrm{Point},W),u}_a\}$ are $\eps$-self-consistent between
	registers $\mathsf{A}$ and $\mathsf{B}$ in $\ket{\hat{\psi}}$
	(item~\ref{enu:qld-cons} of Lemma~\ref{lem:qld-win-implications}), the
	Pauli projections $\{\tau^{W,u}_a\}$ are perfectly self-consistent
	between registers $\mathsf{A}'$ and $\mathsf{A}''$ on the EPR ancillas,
	and the data processing inequality (Lemma~\ref{fact:data-processing})
	combines the two through the convolution form of
	Definition~\ref{def:expanded-point-measurement}. For item~2:
	by~\eqref{eq:lc-22}, each side expands as
	$\frac14 \big( \Id + (-1)^b \hat{X}^{r_X}(u_X) + (-1)^{b'} \hat{Z}^{r_Z}(u_Z) + (-1)^{b+b'} \hat{X}^{r_X}(u_X) \hat{Z}^{r_Z}(u_Z) \big)$
	respectively its reverse-ordered analogue, so it suffices to commute
	the two expanded observables. By~\eqref{eq:lc-23},
	$\hat{X}^{r_X}(u_X)\hat{Z}^{r_Z}(u_Z) = X^{r_X}(u_X) Z^{r_Z}(u_Z) \ot \tau^X(r_X \cdot \mathrm{ind}_m(u_X)) \tau^Z(r_Z \cdot \mathrm{ind}_m(u_Z))$;
	the strategy observables approximately commute or anticommute with sign
	$(-1)^{\gamma(\omega)}$ and error $\sqrt{\eps}$
	(\eqref{eq:pts-obs-commutation} of
	Lemma~\ref{lem:qld-win-implications-obs}), while the generalized Pauli
	observables exactly commute or anticommute with the same sign
	(see Lemma~\ref{lem:twisted-commutation}). The two signs cancel, so
	$\hat{X}^{r_X}(u_X)\hat{Z}^{r_Z}(u_Z) \approx_{\sqrt{\eps}} \hat{Z}^{r_Z}(u_Z)\hat{X}^{r_X}(u_X)$
	on average over uniformly random $\omega$, giving the approximate
	commutation with $\delta_{\mathrm{Anticom}} = \sqrt{\eps}$. The
	symmetric equivalents follow from
	Lemma~\ref{lem:symmetric-equivalents-transfer}.
\end{proof}

\begin{definition}[Expanded line measurements]
	\label{def:expanded-line-measurement}
	% In the source this measurement is defined inside the proof of the lemma
	% labelled lem:qld-comm-line-cons (lines 552-556); it is recorded as a
	% standalone definition because Chapter chap:qpbt-combining uses the
	% concrete measurement, not merely its existence.
	\uses{lem:projective-strategy-setup, def:expanded-state, def:expanded-observables, def:ideg-deg-polynomials, def:low-degree-encoding, lem:pauli-observable-expansion, def:line, def:pauli-win-predicate}
	For $W \in \{X,Z\}$, a line $\ell$, and $f \in \deg_{md}(\ell)$,
	define
	\[
		\hat{M}^{(\mathrm{Line},W),\ell}_f
		\,=\, \sum_{\substack{f',f'' \in \deg_{md}(\ell) \\ f' + f'' = f}}
		M^{(\mathrm{Line},W),\ell}_{f'} \ot \tau^{W,\ell}_{f''}\,,
	\]
	where the projector $\tau^{W,\ell}_{f''}$, acting on register
	$\mathsf{A}'$ (or on $\mathsf{A}''$, following the register conventions
	of Definition~\ref{def:expanded-observables}), is the outcome-$f''$
	projector of the following projective measurement: measure all $M$
	qudits in the basis $W$, obtaining an outcome $h \in \F_q^{M}$, and
	return the polynomial $f''$ obtained by substituting the
	parameterization of $\ell$ into the low-degree encoding
	$g_h \colon \F_q^m \to \F_q$
	(see Definition~\ref{def:low-degree-encoding}): writing $(u_0,v)$ for
	the base point and direction of $\ell$, the outcome $f''$ is
	$g_h(u_0 + tv)$, expanded as a polynomial in the variable $t$ from the
	coefficient representation of $g_h$. Since $g_h$ has individual degree
	at most $d$, hence total degree at most $md$, the substituted
	polynomial has degree at most $md$ in $t$, so
	$f'' \in \deg_{md}(\ell)$. The strategy measurement
	$\{M^{(\mathrm{Line},W),\ell}_{f'}\}$ is indexed by the answers of the
	prescribed form (Definition~\ref{def:pauli-win-predicate}): for a
	diagonal line these are the elements of $\deg_{md}(\ell)$, and for an
	axis-parallel line the elements of $\deg_d(\ell)$, regarded as a
	measurement with outcomes in $\deg_{md}(\ell)$ through the inclusion
	$\deg_d(\ell) \subseteq \deg_{md}(\ell)$ of
	Definition~\ref{def:ideg-deg-polynomials}, with the zero operator
	assigned to outcomes outside $\deg_d(\ell)$. For an axis-parallel
	$\ell$, with direction $e_i$, the substituted polynomial also has
	degree at most $d$ in $t$ --- its degree is the individual degree of
	$g_h$ in the $i$-th variable --- so
	$\hat{M}^{(\mathrm{Line},W),\ell}_f = 0$ unless $f \in \deg_d(\ell)$.
	In all cases
	$\{\hat{M}^{(\mathrm{Line},W),\ell}_f\}_{f \in \deg_{md}(\ell)}$ is a
	projective measurement.
\end{definition}

\begin{lemma}[Consistency of the expanded line measurements]
	\label{lem:qld-comm-line-cons}
	\uses{def:expanded-point-measurement, def:expanded-line-measurement, def:expanded-state, def:line-point-dist, def:approx-question-indexed-operators, def:povm-distance, def:ideg-deg-polynomials, def:symmetric-equivalents}
	There exists a function $\delta_{\mathrm{Line}}(\eps) = \poly(\eps)$
	such that the following holds, with respect to the state
	$\ket{\hat{\psi}}$ of Definition~\ref{def:expanded-state}. For all
	$W \in \{X,Z\}$ and lines $\ell$ there exists a projective measurement
	$\{\hat{M}^{(\mathrm{Line},W),\ell}_f\}_{f \in \deg_{md}(\ell)}$,
	acting on registers $\mathsf{A}\mathsf{A}'$ or
	$\mathsf{B}\mathsf{A}''$, such that:
	\begin{enumerate}
		\item \label{enu:qld-comm-line-self-cons}
		(\textbf{Self-consistency}) For all $W \in \{X,Z\}$ and
		$r \in \F_q$ (the relation does not depend on $r$), on average over
		a uniformly random pair $(\ell,u)$ drawn from the line-point
		distribution (Definition~\ref{def:line-point-dist}),
		\[
			(\hat{M}^{(\mathrm{Line},W),\ell}_f)_{\mathsf{A}\mathsf{A}'}
			\,\approx_{\delta_{\mathrm{Line}}}\,
			(\hat{M}^{(\mathrm{Line},W),\ell}_f)_{\mathsf{B}\mathsf{A}''}\,,
		\]
		where the answer summation is over $f \in \deg_{md}(\ell)$.
		\item \label{eq:qld-comm-line-pt-cons}
		(\textbf{Consistency with points measurements I}) On average over
		a uniformly random pair $(\ell,u)$ drawn from the line-point
		distribution,
		\[
			(\hat{M}^{(\mathrm{Line},W),\ell}_f)_{\mathsf{A}\mathsf{A}'}
			\,\approx_{\delta_{\mathrm{Line}}}\,
			(\hat{M}^{(\mathrm{Line},W),\ell}_f)_{\mathsf{A}\mathsf{A}'}
			\ot (\hat{M}^{(\mathrm{Point},W),u}_{f(u)})_{\mathsf{B}\mathsf{A}''}\,,
		\]
		where the answer summation is over $f \in \deg_{md}(\ell)$.
		\item \label{eq:qld-comm-line-pt-cons2}
		(\textbf{Consistency with points measurements II}) On average over
		a uniformly random pair $(\ell,u)$ drawn from the line-point
		distribution,
		\[
			(\hat{M}^{(\mathrm{Line},W),\ell}_{[\mathrm{eval}_u(\cdot)=a]})_{\mathsf{A}\mathsf{A}'}
			\,\approx_{\delta_{\mathrm{Line}}}\,
			(\hat{M}^{(\mathrm{Point},W),u}_a)_{\mathsf{B}\mathsf{A}''}\,,
		\]
		where
		$\hat{M}^{(\mathrm{Line},W),\ell}_{[\mathrm{eval}_u(\cdot)=a]} = \sum_{f \in \deg_{md}(\ell) \colon f(u)=a} \hat{M}^{(\mathrm{Line},W),\ell}_f$
		and the answer summation is over $a \in \F_q$.
	\end{enumerate}
	Furthermore, all symmetric equivalents of these approximations also
	hold. The proof exhibits the measurements of
	Definition~\ref{def:expanded-line-measurement} as witnesses, with
	item~\ref{eq:qld-comm-line-pt-cons} holding with error $\eps$ and
	items~\ref{enu:qld-comm-line-self-cons}
	and~\ref{eq:qld-comm-line-pt-cons2} with
	$\delta_{\mathrm{Line}}(\eps) = \sqrt{\eps}$.
\end{lemma}
\begin{proof}
	\uses{def:expanded-line-measurement, def:expanded-point-measurement, lem:qld-win-implications, fact:data-processing, fact:add-a-proj2, fact:agreement, def:line-point-dist, def:ideg-deg-polynomials, def:pauli-win-predicate, lem:symmetric-equivalents-transfer}
	Take $\{\hat{M}^{(\mathrm{Line},W),\ell}_f\}$ to be the measurement of
	Definition~\ref{def:expanded-line-measurement}.
	Item~\ref{enu:qld-comm-line-self-cons} follows from the fact that the
	original line measurements $\{M^{(\mathrm{Line},W),\ell}_f\}$ are
	$\eps$-self-consistent between registers $\mathsf{A}$ and $\mathsf{B}$
	in $\ket{\hat{\psi}}$ (item~\ref{enu:qld-cons} of
	Lemma~\ref{lem:qld-win-implications}), the Pauli projections
	$\{\tau^{W,\ell}_{f''}\}$ are perfectly self-consistent between
	registers $\mathsf{A}'$ and $\mathsf{A}''$, and the data processing
	inequality (Lemma~\ref{fact:data-processing}). For
	item~\ref{eq:qld-comm-line-pt-cons}: for the outcomes $f'$ that
	evaluate at $u$, insert the projection
	$M^{(\mathrm{Line},W),\ell}_{[\mathrm{eval}_u(\cdot)=f'(u)]}$ by
	projectivity, replace it by
	$(M^{(\mathrm{Point},W),u}_{f'(u)})_{\mathsf{B}\mathsf{A}''}$ using the
	line-point consistency (item~\ref{enu:qld-ld} of
	Lemma~\ref{lem:qld-win-implications}) together with
	Lemma~\ref{fact:add-a-proj2}, adjoin the exactly consistent projector
	$\tau^{W,u}_{f''(u)}$ on register $\mathsf{A}''$, and recombine by the
	convolution form of Definition~\ref{def:expanded-point-measurement}.
	Outcomes without an evaluation at $u$ occur only when $\ell$ has zero
	direction, where $\tau^{W,\ell}_{f''}$ vanishes unless $f''$ is a
	constant polynomial, so that $f = f' + f''$ evaluates at $u$ exactly
	when $f'$ does; for these outcomes the right-hand side of
	item~\ref{eq:qld-comm-line-pt-cons} carries the zero operator
	(Definition~\ref{def:ideg-deg-polynomials}), and their total mass
	satisfies
	\[
		\E_{(\ell,u)} \sum_{f' \text{ without an evaluation at } u}
		\bra{\hat{\psi}}\, M^{(\mathrm{Line},W),\ell}_{f'} \ot \Id\,
		\ket{\hat{\psi}} \,\leq\, O(\eps)\,,
	\]
	because the left-hand side is exactly the mass of the class
	$M^{(\mathrm{Line},W),\ell}_{[\mathrm{eval}_u(\cdot)=\bot]}$ of the
	completed family in item~\ref{enu:qld-ld} of
	Lemma~\ref{lem:qld-win-implications} --- its value on
	$\ket{\hat{\psi}}$ coincides with its value on $\ket{\psi}$, the
	operator acting trivially on the ancilla registers --- whose partner
	class $M^{(\mathrm{Point},W),u}_{\bot}$ is the zero operator, so that
	the whole of this mass is disagreement mass of that consistency
	statement.
	This gives item~\ref{eq:qld-comm-line-pt-cons} with error $\eps$. For
	item~\ref{eq:qld-comm-line-pt-cons2}, the data processing inequality
	does not apply directly, since the bound of
	item~\ref{eq:qld-comm-line-pt-cons} is stated with $\approx$ rather
	than $\simeq$ and its right-hand side is not of the form
	$\Id \ot B$; instead, a direct computation from
	item~\ref{eq:qld-comm-line-pt-cons} together with projectivity gives
	% The source ends this display with the exact upper bound eps; the bound
	% inherited from item 2 carries the hidden universal constant of
	% def:povm-distance, which is retained inside the O(.).
	\[
		\E_{(\ell,u)} \sum_{a \in \F_q}
		\bra{\hat{\psi}}\, \hat{M}^{(\mathrm{Line},W),\ell}_{[\mathrm{eval}_u(\cdot)=a]}
		\ot \big( \Id - \hat{M}^{(\mathrm{Point},W),u}_a \big) \ket{\hat{\psi}}
		\,\leq\, O(\eps)\,,
	\]
	while the mass bound of the previous paragraph controls the outcomes
	omitted from every class,
	$\E_{(\ell,u)} \bra{\hat{\psi}} \big( \Id - \sum_a
	\hat{M}^{(\mathrm{Line},W),\ell}_{[\mathrm{eval}_u(\cdot)=a]} \big)
	\ot \Id \ket{\hat{\psi}} \leq O(\eps)$; both terms are needed because
	for a zero-direction $\ell$ the classes
	$[\mathrm{eval}_u(\cdot)=a]$ need not exhaust $\deg_{md}(\ell)$.
	An elementary bound for projective sub-measurements $\{A_i\}$,
	$\{B_i\}$ and a state $\ket{\varphi}$, namely
	$\sum_i \bra{\varphi}(A_i - B_i)\ket{\varphi} \leq 2 \big( \sum_i \big\|(A_i - B_i)\ket{\varphi}\big\|^2 \big)^{1/2}$,
	upgrades these two bounds to
	$\Id \approx_{\sqrt{\eps}} \sum_a \hat{M}^{(\mathrm{Line},W),\ell}_{[\mathrm{eval}_u(\cdot)=a]} \ot \hat{M}^{(\mathrm{Point},W),u}_a$.
	A second elementary fact about projective measurements --- if
	$\Id \approx_\delta \sum_i A^x_i \ot B^x_i$ for families of projective
	measurements $\{A^x_i\}$ and $\{B^x_i\}$ acting on separate registers
	of a bipartite state, then $A^x_i \ot \Id \simeq_\delta \Id \ot B^x_i$,
	and hence $A^x_i \ot \Id \approx_\delta \Id \ot B^x_i$ by item~1 of
	Lemma~\ref{fact:agreement} --- is then applied to the two families
	completed as in Definition~\ref{def:ideg-deg-polynomials}: the
	evaluation classes of the expanded line measurement together with
	$\hat{M}^{(\mathrm{Line},W),\ell}_{[\mathrm{eval}_u(\cdot)=\bot]}$,
	and the expanded point measurement together with
	$\hat{M}^{(\mathrm{Point},W),u}_{\bot} = 0$. Both completed families
	are projective measurements, and the completion leaves the sum in the
	hypothesis unchanged, since the term for the outcome $\bot$ vanishes.
	Restricted to the outcomes in $\F_q$, the conclusion yields
	item~\ref{eq:qld-comm-line-pt-cons2} with error $\sqrt{\eps}$, so that
	$\delta_{\mathrm{Line}}(\eps) = \sqrt{\eps}$. The symmetric
	equivalents follow from
	Lemma~\ref{lem:symmetric-equivalents-transfer}.
\end{proof}

The expanded point and line measurements constructed in this section, being
approximately consistent and approximately commuting, are combined into a
single strategy for the classical low individual degree test in
Chapter~\ref{chap:qpbt-combining}.
